% Вставляется в main.tex до leds.tex
\section{Язык Lua}
\subsection{Введение}
Lua — лёгкий и встраиваемый язык программирования, созданный в начале 1990‑х для расширения приложений скриптами. Он компактен, прост и быстро работает, поэтому широко используется в встраиваемых системах, играх и робототехнике.

\subsection{Толковый словарь}
Нажмите на термин в тексте, чтобы перейти к его определению.
\begin{description}
  \item[\hypertarget{term:number}{Число}] вещественное значение (обычно двойной точности), например \verb|3.14| или \verb|10|; доступна математика через \verb|math|.
  \item[\hypertarget{term:string}{Строка}] последовательность символов в кавычках, например \verb|"Lua"|; операции: конкатенация \verb|..|, длина \verb|#s|.
  \item[\hypertarget{term:boolean}{Булево}] логическое значение \verb|true| или \verb|false|; используется в условиях.
  \item[\hypertarget{term:nil}{nil}] «пусто»/«нет значения»; удаляет поля таблицы, означает отсутствие данных.
  \item[\hypertarget{term:table}{Таблица}] универсальная структура (массив и словарь одновременно); хранит позиции \verb|t[1]| и поля \verb|t.x|.
  \item[\hypertarget{term:function}{Функция}] блок кода, вызываемый по имени, принимает аргументы и возвращает значения.
  \item[\hypertarget{term:type}{type()}] функция, возвращающая тип значения: \verb|type(10) -> "number"|.
  \item[\hypertarget{term:scope}{Область видимости}] \verb|local| делает переменную видимой только в текущем блоке; без \verb|local| — глобальная.
  \item[\hypertarget{term:operator}{Оператор}] символ или слово, выполняющее действие: \verb|+ - * / % ^ and or not|.
  \item[\hypertarget{term:expression}{Выражение}] комбинация значений и операторов, дающая новое значение.
  \item[\hypertarget{term:condition}{Условие}] логическое выражение, определяющее выбор ветки выполнения.
  \item[\hypertarget{term:loop}{Цикл}] повторение действий: \verb|for|, \verb|while|, \verb|repeat until|.
  \item[\hypertarget{term:ipairs}{ipairs/pairs}] итераторы по таблицам: \verb|ipairs| для массивов; \verb|pairs| для словарей (порядок не гарантирован).
  \item[\hypertarget{term:truthiness}{Истинность}] в Lua ложными считаются только \verb|false| и \verb|nil|.
  \item[\hypertarget{term:unpack}{unpack}] распаковывает таблицу в список аргументов: \verb|unpack({1,0,0}) -> 1,0,0|.
\end{description}

\subsection{История и назначение}
- Появился как средство «подключаемой логики» поверх приложений на C/C++.
- Минимальные требования к памяти, переносимость и простой синтаксис.
- Табличные структуры данных (\verb|table|) и функции первого класса делают язык выразительным при небольшом ядре.

\subsection{Lua в нашем курсе}
- Мы используем Lua, чтобы программировать поведение квадрокоптера «Геоскан Пионер».
- Через модуль \verb|Ledbar| управляем светодиодной индикацией.
- Через \verb|Timer| создаём неблокирующие анимации.
- Через \verb|callback(event)| реагируем на системные события (посадка, напряжение, достижение точки).

\subsection{Базовые конструкции Lua}
Ниже — краткий пример базовых элементов: переменная, строка и таблица. Циклы и условия разбираются далее.
\begin{minted}{lua}
-- числовая переменная
local x = 10
-- строковая переменная
local s = "Lua"
-- таблица с позиционными элементами и именованным полем
local t = {1, 2, 3}
t.mode = "auto"
-- выводим значения и длину массива
print(x, s, #t)  -- 10  Lua  3
\end{minted}

\subsection{Переменные}
Переменная — это имя, за которым «закреплено» значение. В Lua есть базовые типы: \hyperlink{term:number}{число}, \hyperlink{term:string}{строка}, \hyperlink{term:boolean}{булево}, \hyperlink{term:table}{таблица}, \hyperlink{term:function}{функция}, \hyperlink{term:nil}{nil}. Узнать тип можно через \hyperlink{term:type}{type()}. \hyperlink{term:scope}{Локальная} переменная объявляется через \verb|local| и видна внутри текущего блока; \textbf{глобальная} — без \verb|local| (видна во всём скрипте).

Число — это величина для арифметики (например, \verb|3.14| или \verb|10|). Строка — текст в кавычках (например, \verb|"Pioneer"|), над ней выполняют конкатенацию \verb|..| и получают длину \verb|#s|. Булево — \verb|true|/\verb|false|, используется в \hyperlink{term:condition}{условиях}. Таблица — универсальная структура: массив \verb|t[1]| и словарь \verb|t.x| одновременно. Функция — именуемый блок кода, например \verb|local function f(x) return x*2 end|.

\begin{minted}{lua}
-- локальная числовая переменная
local x = 10
-- глобальная строковая переменная
name = "Pioneer"
-- локальная булева переменная
local ready = true
-- таблица с позиционными элементами и именованным полем
local t = {1, 2, 3, speed = 3.5}
-- множественное присваивание
x, y = 1, 2
-- обмен значениями
x, y = y, x
-- выясняем типы значений
print(type(x), type(name), type(ready), type(t))  -- number string boolean table
\end{minted}
\begin{minted}{lua}
-- словарь (именованные поля)
local p = {x = 1, y = 2}
-- массив (позиционные элементы)
local a = {10, 20, 30}
-- добавление нового поля
p.z = 3
-- добавление нового элемента массива
a[4] = 40
\end{minted}
\paragraph{Пояснения.}
- Таблица — единственная сложная структура в Lua; через неё строят списки, словари и объекты.
- \verb|nil| при присваивании полю таблицы удаляет это поле: \verb|t.speed = nil|.
- Глобальные переменные удобно использовать только для настроек и состояний, иначе код становится хрупким; предпочитайте \verb|local|.

\subsection*{Задачи для самостоятельного решения}
\begin{enumerate}
  \item Объявите локальную переменную \verb|x| со значением \verb|42|.
  \item Присвойте строку \verb|"Pioneer"| глобальной переменной \verb|name|.
  \item Поменяйте местами значения \verb|a| и \verb|b| одной строкой.
  \item Создайте таблицу \verb|t| с полем \verb|speed=3.5| и элементами \verb|1,2,3|.
  \item Добавьте в таблицу \verb|t| поле \verb|mode| со значением \verb|"auto"|.
  \item Объявите \verb|local flag = false| и затем присвойте глобальной \verb|flag| значение \verb|true|.
  \item Создайте пустую таблицу \verb|buf| и добавьте в неё числа \verb|5| и \verb|7|.
  \item Удалите поле \verb|speed| из таблицы \verb|t|, присвоив \verb|nil|.
  \item Присвойте одновременно значения \verb|r=1|, \verb|g=0|, \verb|b=0|.
  \item Создайте таблицу \verb|color| вида \verb|{r=1,g=0,b=0}|.
  \item Извлеките третий элемент массива \verb|a| в переменную \verb|third|.
  \item Присвойте \verb|local s = "Lua"| и \verb|local n = 3.14|.
  \item Создайте локальную функцию \verb|id|, возвращающую свой аргумент.
  \item Объявите \verb|local config = nil| и затем присвойте \verb|config = {enabled=true}|.
  \item Создайте таблицу \verb|point| и добавьте в неё \verb|x=10| и \verb|y=20|.
\end{enumerate}

\subsection{Арифметические операции}
Операции: \verb|+ - * / % ^| (плюс, минус, умножение, деление, остаток, возведение в степень). Сравнения: \verb|< <= > >= == ~=|. В Lua 5.1 нет операторов целочисленного деления \verb|//| и возведения \verb|**|; используйте \verb|math.floor(a/b)| и \verb|a ^ b| соответственно. Все числа — \hyperlink{term:number}{вещественные}, поэтому деление отдаёт дробные значения; для округлений применяются \verb|math.floor|, \verb|math.ceil| и \verb|math.abs|.

Приоритет: сначала выполняются \verb|^|, затем \verb|* / %|, далее \verb|+ -|. Скобки меняют порядок вычислений. Остаток \verb|%| определён как «остаток от деления», полезен для задач чётности и циклических индексов.
\begin{minted}{lua}
-- приоритет: умножение выполняется раньше сложения
local a = 2 + 3 * 4
-- остаток от деления (modulo)
local r = 5 % 2
-- возведение 2 в степень 3
local p = 2 ^ 3
-- сравнение: равно?
local eq = (a == 14)
-- сравнение: больше?
local gt = (a > 10)
-- безопасное «целочисленное» деление с округлением вниз
local q = math.floor(17 / 4)  -- 4
-- евклидова длина вектора
local len = math.sqrt(3*3 + 4*4)  -- 5
\end{minted}
\paragraph{Пояснения.}
- \verb|math.sqrt| возвращает корень квадратный; \verb|math.pi| — число $\pi$.
- \verb|math.max| и \verb|math.min| ограничивают значения: \verb|v = math.max(0, math.min(1, v))|.
- Деление на ноль следует избегать: проверяйте знаменатель перед делением.
\subsection*{Задачи для самостоятельного решения}
\begin{enumerate}
  \item Вычислите выражение \verb|3 + 4 * 2 - 5|.
  \item Найдите остаток \verb|17 % 4|.
  \item Запишите \verb|3 ^ 4|.
  \item Вычислите среднее \verb|(x + y) / 2|.
  \item Запишите \verb|math.abs(x2 - x1)|.
  \item Вычислите \verb|math.floor(n / 3)|.
  \item Запишите условие «\verb|n| чётное» с \verb|%|.
  \item Вычислите \verb|math.sqrt(x*x + y*y)| (евклидова длина).
  \item Ограничьте \verb|v| в диапазоне \verb|[0,1]|: \verb|v = math.max(0, math.min(1, v))|.
  \item Переведите градусы в радианы: \verb|rad = deg * math.pi / 180|.
  \item Вычислите площадь прямоугольника \verb|S = w * h|.
  \item Запишите среднее трёх чисел \verb|(a+b+c)/3|.
  \item Реализуйте линейную интерполяцию \verb|lerp = a + t*(b-a)|.
  \item Посчитайте сумму массива \verb|a|: примените цикл \verb|for|.
  \item Найдите максимум массива \verb|a|: пройдитесь и сравнивайте.
\end{enumerate}

\subsection{Логические операции}
Логика: \verb|and| (И), \verb|or| (ИЛИ), \verb|not| (НЕ). В Lua ложными считаются только \verb|false| и \verb|nil| (\hyperlink{term:truthiness}{истинность}). Операторы \verb|and| и \verb|or| обладают \textbf{коротким замыканием}: правая часть вычисляется только при необходимости. Сравнения работают для чисел и строк (лексикографически).
\begin{minted}{lua}
-- логическое И и НЕ
local ok = true and not false
-- «тернарный» идиом в Lua через and/or
local sel = (gt and 1 or 0)
-- композиция условий «в диапазоне»
local within = (a >= 0) and (a <= 100)
-- проверка непустой строки
local nonempty = (name ~= nil) and (name ~= "")
-- исключающее ИЛИ для булевых
local xor = (p and not q) or (q and not p)
\end{minted}
\paragraph{Пояснения.}
- Идиома \verb|(cond and A or B)| выбирает \verb|A| при истинном \verb|cond| и \verb|B| иначе.
- \verb|nil| часто означает «нет данных», поэтому проверяйте значение перед доступом к полям.
- Строки сравниваются лексикографически: \verb|"a" < "b"| истинно.
\subsection*{Задачи для самостоятельного решения}
\begin{enumerate}
  \item Составьте условие «\verb|x| в диапазоне \verb|[0,1]|».
  \item Верните \verb|1| если \verb|b==true|, иначе \verb|0| с \verb|and/or|.
  \item Составьте проверку «строка \verb|s| не пустая» (\verb|s~=""|).
  \item Сравните \verb|"a"| и \verb|"b"|: проверка «меньше» (\verb|<|).
  \item Составьте выражение «\verb|x>0| и \verb|y~=0|».
  \item Запишите де Моргана: \verb|not(a and b)| эквивалент \verb|(not a) or (not b)|.
  \item Реализуйте «исключающее ИЛИ» для \verb|p| и \verb|q|.
  \item Преобразуйте число \verb|n| в булево: \verb|n~=0|.
  \item Составьте логическое выражение: \verb|t.enabled == true or t.mode=="auto"|.
  \item Верните «сигнал опасности», если \verb|voltage<3.5| или \verb|temp>80|.
  \item Напишите предикат «индекс чётный»: \verb|i % 2 == 0|.
  \item Составьте выражение, выбирающее максимум из \verb|a| и \verb|b| через \verb|and/or|.
  \item Используя \verb|not|, инвертируйте булеву переменную \verb|on|.
  \item Составьте условие «\verb|t| не \verb|nil| и \verb|t.value| существует».
  \item Напишите логическую функцию «входит ли символ \verb|c| в набор \verb|{"r","g","b"}|».
\end{enumerate}

\subsection{Условия}
Ветвления выполняются через \verb|if ... then ... elseif ... else ... end|. Истинными считаются любые значения, кроме \verb|false| и \verb|nil|. Шаблоны: «охранные условия» (ранний выход), таблица переходов (эмуляция конвейера состояний), каскад \verb|elseif| для диапазонов.
\begin{minted}{lua}
-- классификация числа x
if x > 0 then
  state = "positive"
elseif x == 0 then
  state = "zero"
else
  state = "negative"
end
\end{minted}
\begin{minted}{lua}
-- безопасно работаем с полями, если таблица и поле существуют
if t and t.value then
  result = t.value * 2
end
\end{minted}
\paragraph{Пояснения.}
- «Охранное» условие позволяет быстро завершить ветку при неверных входных данных.
- Вложенные условия стоит разбивать на логические блоки для читаемости.
- Для выбора из фиксированного набора значений удобно использовать таблицу функций: \verb|actions[state]()|.
\subsection*{Задачи для самостоятельного решения}
\begin{enumerate}
  \item Напишите условие, присваивающее \verb|sign="+"| если \verb|x>0|, иначе \verb|sign="-"|.
  \item Присвойте \verb|parity="even"| при чётном \verb|n|, иначе \verb|"odd"|.
  \item Составьте ветвление: \verb|temp<0| → \verb|"cold"|, \verb|0<=temp<20| → \verb|"cool"|, иначе \verb|"warm"|.
  \item Присвойте \verb|level="low"| если \verb|v<3.7|, \verb|"ok"| если \verb|v>=3.7|.
  \item Напишите условие: если \verb|config.enabled| истинно, то \verb|mode="run"|, иначе \verb|"stop"|.
  \item Реализуйте проверку: если \verb|t| не \verb|nil| и \verb|t.x| существует, удвойте \verb|t.x|.
  \item Напишите ветвление, присваивающее цвет: \verb|r=1,g=0,b=0| при \verb|state=="RED"|, \verb|r=1,g=1,b=0| при \verb|"YELLOW"|, иначе \verb|r=0,g=1,b=0|.
  \item Составьте условие: если \verb|s==""|, то \verb|empty=true|, иначе \verb|false|.
  \item Реализуйте защиту от деления на ноль: если \verb|b~=0|, то \verb|q=a/b|.
  \item Напишите проверку «строка \verb|name| начинается с \verb|"P"|» (используйте \verb|string.sub|).
  \item Если \verb|voltage<3.5|, установите \verb|alarm=true|.
  \item Составьте ветвление: если \verb|t.mode=="manual"| и \verb|t.enabled|, то \verb|run=true|.
  \item Если \verb|count==0|, присвойте \verb|count=1|, иначе увеличьте на \verb|1|.
  \item Напишите условие, которое инвертирует булеву переменную \verb|on|.
  \item Реализуйте тройное ветвление по \verb|speed|: \verb|slow|, \verb|medium|, \verb|fast|.
\end{enumerate}

\subsection{Циклы}
Циклы: числовой \verb|for| (счётчик от \verb|start| до \verb|finish| с шагом), \verb|while| (пока условие истинно), \verb|repeat ... until| (выполнить хотя бы один раз). Для таблиц: \hyperlink{term:ipairs}{ipairs} (массив) и \verb|pairs| (словарь). \verb|break| прерывает цикл; \verb|continue| отсутствует (можно имитировать через \verb|goto|, но лучше структурировать условие).
\begin{minted}{lua}
-- суммирование 1..10
local sum = 0
for i = 1, 10 do
  sum = sum + i
end
\end{minted}
\begin{minted}{lua}
-- увеличиваем x до 5
local x = 0
while x < 5 do
  x = x + 1
end
-- уменьшаем x до 0 минимум один раз
repeat
  x = x - 1
until x == 0
\end{minted}
\begin{minted}{lua}
-- умножаем элементы массива на 2
local a = {10, 20, 30}
for i, v in ipairs(a) do
  a[i] = v * 2
end
-- увеличиваем значения словаря на 1
local t = {x=1, y=2}
for k, v in pairs(t) do
  t[k] = v + 1
end
\end{minted}
\paragraph{Пояснения.}
- \verb|ipairs| проходит по индексам 1..N без пропусков; \verb|pairs| проходит по всем полям словаря в произвольном порядке.
- В числовом \verb|for| границы включительно: \verb|for i=1,3 do| даёт \verb|1,2,3|.
- Для обхода диапазона индексов, используемых в светодиодной линейке, применяйте \verb|for i=0,24 do|.
\subsection*{Задачи для самостоятельного решения}
\begin{enumerate}
  \item Посчитайте сумму чисел от \verb|1| до \verb|100|.
  \item Найдите максимум в массиве \verb|a|.
  \item Подсчитайте количество чётных элементов в массиве \verb|a|.
  \item Постройте новый массив \verb|b| из квадратов элементов \verb|a|.
  \item Разверните массив \verb|a| (создайте \verb|rev|).
  \item Удалите из массива \verb|a| все нули.
  \item Сформируйте массив \verb|seq| из чисел \verb|1..20|.
  \item Вычислите факториал \verb|n| с помощью \verb|while|.
  \item Реализуйте цикл \verb|repeat ... until|, который увеличивает \verb|x| до \verb|>=10|.
  \item Скопируйте словарь \verb|t| в \verb|copy|, увеличив каждое значение на \verb|1|.
  \item Объедините два массива \verb|a| и \verb|b| в \verb|c|.
  \item Посчитайте, сколько значений \verb|true| в массиве булевых \verb|flags|.
  \item Создайте таблицу \verb|squares| с полями \verb|["1"]=1, ["2"]=4, ..., ["10"]=100|.
  \item Пройдитесь по индексам \verb|0..24| и заполните таблицу \verb|idx| значениями \verb|i|.
  \item Постройте массив \verb|fibs| из первых \verb|10| чисел Фибоначчи.
\end{enumerate}

\subsection{Работа с модулями платформы}
- \verb|Ledbar.new(N)| создаёт объект управления светодиодами (используем \verb|N=25|).
- \verb|leds:set(i, r, g, b)| задаёт цвет по индексу \verb|i| (\verb|r,g,b| в \verb|[0.0,1.0]|).
- \verb|Timer.new(period, fn)| создаёт периодический таймер; \verb|Timer.callLater(delay, fn)| — отложенный вызов.
- \verb|function callback(event)| — обработчик системных событий.

\subsection{Примеры, используемые далее}
\begin{minted}{lua}
local unpack = table.unpack
local ledNumber = 25
local leds = Ledbar.new(ledNumber)
local function changeColor(col)
  for i=0, ledNumber-1 do
    leds:set(i, unpack(col))
  end
end
changeColor({0,1,0})
\end{minted}
\begin{minted}{lua}
local ledNumber = 25
local leds = Ledbar.new(ledNumber)
for i=0, ledNumber-1 do leds:set(i, 0, 0, 0) end
Timer.callLater(3, function ()
  for i=0, ledNumber-1 do leds:set(i, 0, 0, 1) end
end)
\end{minted}
\begin{minted}{lua}
local ledNumber = 25
local leds = Ledbar.new(ledNumber)
on = false
function updateBlink()
  local r = on and 1 or 0
  for i=0, ledNumber-1 do leds:set(i, r, 0, 0) end
  on = not on
end
blinkTimer = Timer.new(0.5, updateBlink)
blinkTimer:start()
\end{minted}
